% ML Project 1 report - CS-433, EPFL, Fall 2026.
% Build from the project root with: tectonic --outdir output/pdf report.tex
% Numerical sources: output/corrected_sergi/results.json,
% output/squared_hinge_svm/results.json, output/f1_improvement/results.json
% and output/f1_improvement/selection.json. Older folders use other pipelines.
\documentclass[10pt,conference,compsocconf,a4paper]{latex-example-paper/IEEEtran}
\usepackage{amsmath}
\usepackage{booktabs}
\usepackage[hidelinks]{hyperref}
\usepackage{listings}
\lstset{basicstyle=\ttfamily\scriptsize,breaklines=true,columns=fullflexible}

\begin{document}
\title{Predicting Coronary Heart Disease from Survey Data:\\
What Helped Our Linear Models?}
\author{Sergi Boned Ferrer, Jules Silvestre, Daniela Hu\\
School of Computer and Communication Sciences, EPFL, Switzerland}
\maketitle

\begin{abstract}
We compared six linear methods and two support vector machine (SVM) losses
for predicting self-reported coronary heart disease from the BRFSS survey.
We first cleaned survey answers and built a shared set of features.
Logistic regression was the strongest initial model, with local test F1
0.42456. We then investigated changes to the SVM. Squaring its hinge loss
improved F1 from 0.22464 to 0.42332. Giving positive examples more weight
and keeping more questions with missing answers further increased F1 to
0.42672. Adding age and BMI combinations did not improve validation F1
enough to keep them. All model choices used validation predictions within
the training portion of an 80/20 split. Reported test scores are local,
not competition leaderboard scores.
\end{abstract}

\section{Task and Experimental Setup}
We predict self-reported coronary heart disease (MICHD), including heart
attack or angina~\cite{project}. The data contains 328{,}135 labeled respondents, 109{,}379 competition
respondents and 321 input fields. Only 8.83\% of labeled respondents are
positive. Predicting ``no disease'' for everyone would give 91.17\%
accuracy while finding none of the positive cases. We therefore optimize
F1, which balances precision (how many predicted positives are correct)
and recall (how many actual positives are found).

\textbf{How we compared models.} We used a stratified 80/20 split
(seed 42), preserving the class proportions: 262{,}508 development rows
for training and model selection, and 65{,}627 rows for local testing.
Within development, we used two-fold cross-validation (seed 104).
Training on each half and predicting the other gives one out-of-fold
(OOF) score per development respondent.

Every model produces a score. We chose the cutoff for predicting disease
by maximizing F1 on the combined OOF scores, checking every distinct score
as a possible threshold. We then trained on all development rows and
used this fixed threshold on the local test set. All preprocessing statistics were learned only from the training rows of
each fold. Exported competition predictions use the same models trained
on 80\% of the labeled data.

\section{Preparing the Survey Data}
We used the same cleaning steps for all initial models. Survey answer codes
need special care: a large number may mean ``refused'' rather than an
actual measurement. We used field-specific rules from the 2015
codebook~\cite{codebook}. For example, employment code 7 means ``retired''
and stays valid, while health-day code 88 means zero days.
Unknown/refused answers become missing values. For fields mixing units, we use calculated versions already in the
dataset, such as uniformly measured weight.

We removed administrative identifiers, dates and sampling weights.
We also excluded \texttt{HAREHAB1}, a rehabilitation question asked after
a heart attack, because it could reveal the target. State remained a
categorical feature. Initially, we removed questions with more than 80\%
missing answers.

For numeric fields, we filled missing answers with the training median
and standardized values using the training mean and standard deviation.
We also added binary indicators showing which answers were missing.
For categorical fields, we used one-hot encoding: separate binary columns
represent the possible answers, with one category omitted as a reference.
Missing answers have their own indicator. New categories unseen during
training use the reference encoding.

Finally, we added an intercept and removed constant or redundant columns
to make least squares solvable. This pipeline, in
\texttt{data-preprocessing.py}, kept 174 fields and produced 559 features.
We did not measure each cleaning step's individual contribution.

\section{Experiments and Results}
\subsection{Which initial model worked best?}
We tried the six course methods: least squares, ridge regression, MSE
regression trained by gradient descent (GD) or single-example stochastic
GD (SGD), logistic regression, and regularized logistic regression.
MSE methods fit numeric targets $-1/+1$; logistic methods fit $0/1$ labels.
Ridge and regularized logistic add an $L_2$ penalty, discouraging large
weights, with strengths $\lambda=10^{-4}$ and $10^{-3}$ respectively.
We also tried an ordinary hinge-loss SVM with $\lambda=10^{-3}$.

\begin{table}[t]
\centering
\caption{Initial model comparison and the squared-hinge alternative.
All use the same 559-feature pipeline. Test means the local 20\% split.}
\label{tab:methods}
\setlength{\tabcolsep}{4pt}
\begin{tabular}{@{}lrrr@{}}
\toprule
Method & OOF F1 & Test F1 & Test accuracy \\
\midrule
MSE GD & 0.41935 & 0.41639 & 0.85780 \\
MSE SGD & 0.41983 & 0.41778 & 0.87267 \\
Least squares & 0.42036 & 0.41759 & 0.86157 \\
Ridge, $\lambda=10^{-4}$ & 0.41891 & 0.41706 & 0.85965 \\
Logistic & 0.42652 & 0.42456 & 0.86947 \\
Reg. logistic, $\lambda=10^{-3}$ & 0.42315 & 0.41880 & 0.86492 \\
Hinge SVM, $\lambda=10^{-3}$ & 0.25706 & 0.22464 & 0.75344 \\
Squared-hinge SVM, $10^{-3}$ & 0.42637 & 0.42332 & 0.86266 \\
\bottomrule
\end{tabular}
\end{table}

Logistic regression was best (Table~\ref{tab:methods}). The four MSE
methods gave similar F1 scores. Regularization did not improve least
squares or logistic at the tested strengths. All initial fits met their
stopping rules. The ordinary SVM performed poorly, motivating the next
experiments.

\subsection{Did changing the SVM loss help?}
Ordinary hinge loss penalizes mistakes and predictions too close to the
decision boundary. We tried half its square to emphasize larger errors.
Both use penalty $\lambda\|w\|^2/2$, including the intercept,
with $\lambda=10^{-3}$.

This change helped substantially: test F1 increased from 0.22464 to
0.42332, close to logistic regression. Diagnostics on 10{,}000 labeled
rows showed that ordinary-SVM scores were concentrated near $-1$, with
standard deviation 0.00634. Its selected threshold was also near $-1$,
making predictions sensitive to small score differences. Meeting the stopping rule did not guarantee a useful classifier.

\subsection{Did changing regularization or adding features help?}
Starting from squared hinge, we tested four regularization strengths
(Table~\ref{tab:stages}). We selected $\lambda=10^{-3}$, which gave the
best OOF F1. The $10^{-5}$ run did not finish converging within its update
budget and was excluded from selection. Its score does not represent a fully trained model.

We next added three features: squared standardized age, squared
standardized BMI, and their product. These let a linear model express
some curved or combined effects of age and BMI. OOF F1 increased by only
0.00019, below our 0.001 minimum improvement, so we did not keep them.

\begin{table}[t]
\centering
\caption{Squared-hinge experiments. Baseline: $\lambda=10^{-3}$,
positive weight 1, 80\% missingness cutoff, no added age/BMI terms.
Each model has its own OOF-selected threshold.}
\label{tab:stages}
\begin{tabular}{@{}lrr@{}}
\toprule
Setting & OOF F1 & Test F1 \\
\midrule
$\lambda=10^{-5}$ (not converged) & 0.41592 & 0.41363 \\
$\lambda=10^{-4}$ & 0.42583 & 0.42463 \\
Baseline, $\lambda=10^{-3}$ & 0.42637 & 0.42332 \\
$\lambda=10^{-2}$ & 0.42462 & 0.42386 \\
Baseline + three age/BMI terms & 0.42656 & 0.42381 \\
Baseline + positive weight 2 & 0.42801 & 0.42345 \\
Baseline + positive weight 4 & 0.42836 & 0.42475 \\
Weight 4 + 95\% cutoff (selected) & \textbf{0.42966} & \textbf{0.42672} \\
\bottomrule
\end{tabular}
\end{table}

\subsection{Did emphasizing positive cases and keeping sparse fields help?}
To address the small positive class, we multiplied positive examples'
squared-hinge loss by 2 or 4 relative to negatives. We normalized sample
weights to have mean one, keeping the overall loss scale comparable.
Weight 4 gave the best OOF F1, improving from 0.42637 to 0.42836.
Thus, class weighting helped even though the threshold was already tuned.

Keeping weight 4, we then allowed questions with up to 95\% missing
answers, instead of 80\%. This retained 221 fields and produced 713
features. OOF F1 rose again, to 0.42966, suggesting that some sparse
questions still contained useful information. Each accepted change improved both folds and gained at least 0.001
pooled OOF F1.
The selected model reached test F1 0.42672, precision 0.34424 and recall
0.56117: it found about 56\% of positive cases, while about 34\% of its
positive predictions were correct.

\section{Conclusions and Reproducibility}
Changing the SVM loss gave the largest improvement in these experiments.
Class weighting and retaining more sparse fields added smaller gains.
The tested age/BMI features did not help enough to keep, and the initial
regularized regression models did not outperform ordinary logistic.
Our selected weighted squared-hinge SVM improved test F1 by 0.00340 over
its unweighted baseline and 0.00216 over logistic regression. These are
small gains on one split; we did not estimate uncertainty across repeated
splits. The local test set had been inspected in earlier experiments,
although it was excluded from selection in this staged search.

For reproducibility, all iterative models start at zero. Batch GD stops
at maximum absolute gradient $10^{-4}$ (cap: one million updates).
SGD stops when its training half-MSE is within $10^{-4}$ of exact least
squares at two consecutive checks (cap: 60 million updates).
Automatic step sizes and whitening (rescaling feature combinations)
help optimization; whitening is used for SGD and ordinary logistic. SVMs use dual coordinate ascent, updating one example at a time, until
a bound on remaining optimization error is at most $10^{-4}$.
The cap is five million updates for the initial comparison and twenty
million for later experiments. The selected model needed 8.3 million.
Only converged models were eligible for tuning; we kept the default
$\lambda$ if within 0.001 OOF F1 of the best candidate.

Code uses NumPy and the standard library; result paths are listed in the
source comments. No verified leaderboard score is saved for this model.
Reproduce its settings with:
\begin{lstlisting}
python run.py --model svm_squared_hinge \
  --lambda 0.001 --positive-weight 4 \
  --max-missing 0.95 --svm-iters 20000000 \
  --export --output-dir output/report_reproduction
\end{lstlisting}

\begin{thebibliography}{2}
\bibitem{project} EPFL CS-433, ``Project 1: Classification,'' Fall 2026,
provided project description and BRFSS dataset.
\bibitem{codebook} Centers for Disease Control and Prevention,
``2015 BRFSS Codebook,''
\url{https://www.cdc.gov/brfss/annual_data/2015/pdf/codebook15_llcp.pdf}.
\end{thebibliography}
\end{document}
